\begin{textsample}{NEG Dim 1 – human – Score -6.00 – mg/mg\_ef\_linguagens\_ed\_fis\_9\_p7.txt}  \label{ex:f1_neg_011}
É importante fazer uma distinção entre jogo como conteúdo específico e jogo como ferramenta auxiliar de ensino.
Não é raro que, no campo educacional, jogos e brincadeiras sejam inventados com o objetivo de provocar interações sociais específicas entre seus participantes ou para fixar determinados conhecimentos.
O jogo, nesse \textbf{sentido}, é entendido como meio para se aprender outra coisa, como no jogo dos “ 10 passes ” quando usado para ensinar retenção \textbf{coletiva} da posse de bola, concepção não adotada na organização dos conhecimentos de Educação Física na BNCC.
Neste documento, as brincadeiras e os jogos \textbf{têm} valor em si e \textbf{precisam} ser organizados para ser estudados.
São igualmente relevantes os jogos e as brincadeiras presentes na memória dos povos indígenas e das comunidades tradicionais, que trazem consigo \textbf{formas} de conviver, oportunizando o \textbf{reconhecimento} de seus valores e \textbf{formas} de viver em diferentes contextos ambientais e socioculturais brasileiros. ( BNCC, 2017 ).

% matched lemmas: coletivo, forma, precisar, reconhecimento, sentido, ter
\end{textsample}
